Die Logarithmusfunktion $y(z) = \log (1+z)$ erfüllt die Differentialgleichung
\begin{equation}
(1+z)y'(z) = 1.
\label{buch:differentialgleichungen:501:eqn:dgl}
\end{equation}
Verwenden Sie einen Potenzreihenansatz, um eine Potenzreihe für $\log(1+z)$
zu finden.

\begin{loesung}
Wir schreiben die Potenzreihe als
\[
y(z)
=
\sum_{k=0}^\infty
a_k z^k
\]
und bilden die Ableitung
\[
y'(z)
=
\sum_{k=1}^\infty
a_k k z^{k-1}
=
\sum_{l=0}^\infty
a_{l+1}(l+1)z^l.
\]
Diese Reihe setzen wir jetzt in die Differentialgleichung
\eqref{buch:differentialgleichungen:501:eqn:dgl}
ein und erhalten
\begin{align}
1
&=
(1+z)
\sum_{l=0}^\infty
a_{l+1}(l+1)z^l
\notag
\\
&=
\sum_{l=0}^\infty
a_{l+1}(l+1)z^l
+
\sum_{l=0}^\infty
a_{l+1}(l+1)z^{l+1}
\notag
\\
&=
a_1
+
\sum_{l=1}^\infty
a_{l+1}(l+1)z^l
+
\sum_{l=1}^\infty
a_{l}lz^l
\notag
\\
&=
a_1
+
\sum_{l=1}^\infty
(a_{l+1}(l+1)+a_{l}l)z^l.
\label{buch:differentialgleichungen:501:eqn:reihe}
\end{align}
Durch Koeffizientenvergleich kann man aus
\eqref{buch:differentialgleichungen:501:eqn:reihe}
die Rekursionsbeziehung
\[
a_{l+1} = -\frac{l}{l+1} a_l
\quad\text{mit der Anfangsbedingung}\quad a_1=1
\]
ablesen.
Sie hat die Lösung
\[
a_l
=
(-1)^{l+1}\frac{1}{l} a_1,
\]
denn sie erfüllt einerseits die Anfangsbedingung und andererseits
wegen
\[
a_{l+1}
=
-\frac{l}{l+1}
\biggl(
(-1)^{l+1}\frac{1}{l}
\biggr)
=
(-1)^{l+1+1}
\frac{1}{l}
\frac{l}{l+1}
=
(-1)^{l+2}
\frac{1}{l+1}
\]
die Rekursionsgleichung.
Die gesuchte Differentialgleichung ist daher
\[
\log(1+z)
=
z
-\frac{z^2}{2}
+\frac{z^3}{3}
-\frac{z^4}{4}
+\frac{z^5}{5}
-\frac{z^6}{6}
+\frac{z^7}{7}
-\frac{z^8}{8}
+\frac{z^9}{9}
-\dots
\]
Natürlich kann diese Reihe auch aus der umgeformten Differentialgleichung
\[
y'(z)
=
\frac{1}{1+z}
=
\sum_{k=0}^\infty (-1)^kz^k
=
1-z+z^2-z^3+z^4-z^5+z^6-z^7+z^8-\dots
\]
durch gliedweise Integration der geometrischen Reihe erhalten werden.
\end{loesung}


